% Labels of the draft are printed as \dref{label} (texttt), since this note is compiled on its own.
\providecommand{\dref}[1]{\texttt{\detokenize{#1}}}
\section{Referee 3, Major 5 and the order of limits: prescriptions}\label{sec:ref3presc}
% Scripts: notes/scripts/ref3_presc/  a1_onshell_mom.py (a1.out), a2_rowIII_delta.py (a2.out), a3_R_symbolic.py (a3.out),
% a4_coincidence.py (a4.out), a5_limits.py (a5.out), b1_inst_distr.py (b1.out), c1_order_limits.py (c1.out).
% Sources read: paper/draft/sec_setup.tex, sec_observable.tex, sec_finite.tex (6.4, 6.5, 6.18-6.19, conditions),
% sec_dglap_rc.tex (pole account), app_finite.tex, app_wavefunctions.tex, app_unitarity.tex, app_distributions.tex;
% calc_main_checked/Calculation_main_1_checked.tex (single-gluon dressing, B_3); group_I_checked/Group_I_checked.tex (Row III
% before and after the p+ integral); notes/fix_soft.tex (Steps 8, 9), fix_poles_dglap.tex, scripts/s7_range_ct.py, s8_range_norm.py.

\noindent\textbf{Result in one paragraph.}
(a) In the time-dependent (Gell-Mann--Low) construction the vanishing energy denominators $E_g(k)-E_g(p)$ of $\mf B_3$ acquire
$\pm i\epsilon$; they occur only in $\mf B_3$, which enters the spectrum only through the final-state dressing $\bar\Omega^\dagger$, i.e.\ only in Group~I Row~III ($\mathrm{LO}^\dagger\times\bar{\mf B}_3^\dagger\mf A$) and its complex conjugate. The $i\pi\delta$ part of $\mf B_3$ is anti-Hermitian (checked on the energy shell to $4\times10^{-13}$), so it is a unitary rotation inside the energy shell, allowed by $\Omega^\dagger\Omega=1$ and not fixed by it. In Row~I.III the extra terms are \emph{not} cancelled: the $i\pi\theta$ pieces of the separate terms $N_1,N_2,N_{3a}$ combine \emph{exactly} into a $\theta$-free constant, $\ln\frac{(x-z)^2}{(x-w)^2}\to\ln\frac{(x-z)^2}{(x-w)^2}-i\eta\pi$ in the combined kernel $\mathcal M^{ij}$, Eq.~(\dref{eq:Mij}) ($\eta=+1$ for the out-state dressing, which is the physical one for the observed gluon; $\eta=-1$ for an in-state dressing of the final state; PV is the average). The coefficient of that logarithm is a rotation matrix, $\mathfrak B^{ij}=[\cos2\phi\,\delta^{ij}+\sin2\phi\,\epsilon^{ij}]/[2(\v z-\v w)^2]$, $\phi=\angle(\v x-\v w,\v x-\v z)$. Since the PV term carries $i\times$(real kernel) and the new term a real kernel, the complex conjugate adds rather than cancels. The new term $\Delta N$ is finite only together with its complex conjugate: at $\v z\to\v w$ it has a non-integrable $1/(\v z-\v w)^2$, cut off by the window $|p^+-k^+|>\Lambda$ of $\mf B_3$ into $\ln(k^+/\Lambda)$; that logarithm multiplies $\mathcal A_{\rm LO}^\dagger(f\cdot\bar Q)\mathcal A_{\rm LO}$ ($\bar Q$ the colour charge of the projectile after the collision), i.e.\ it is a forward Coulomb phase, and cancels exactly against the complex conjugate. The remainder, Eq.~\eqref{pr:DN}, is finite and integrable, conserves the total number of gluons, is even in $\v k$ in its symmetrised three-$\rho$ part and odd in its commutator part (opposite to the PV terms), and is not zero in general. \emph{The finite terms of the draft are correct for the principal-value (standing-wave) dressing; for the adiabatic out-state definition of the observed gluon they miss $\Delta N$.} (b) The instantaneous Group~III terms do \emph{not} come from $p^+$ near $\vee$: after the $\v y'$ integral their $p^+$ distribution falls like $(k^+/p^+)^2$, and the bulk comes from $p^+\lesssim k^+\max\{1,1/(|\v k|\ell)\}$, $\ell$ the shortest transverse scale of the integrand; $p^+>P$ contributes $O(k^+/P)$ [$O((k^+/P)^\alpha)$ for a H\"older-$\alpha$ target]. With any smooth separation $f(p^+/\vee)$ the finite parts have the same limit; the only cut-off dependence of the whole result is the universal shift $\ln\vee\to\ln\vee+c_f$ of the BFKL logarithm. They are under control. (c) The leftover $-\frac{\ell_\vee}{2}\frac2\epsilon$ comes from parents faster than $\vee$ in the part of the vertex correction that is not enhanced at small momentum fraction; its coefficient is cut-off-shape dependent ($\ell_\vee\to-\kappa_f k^+/\vee$, $\kappa_f=1,\sqrt\pi,\pi/2$ for three cut-offs). It is removed by defining the effective theory at leading power in $k^+/\vee$ at fixed $\epsilon$; Section~\ref{pr:c} gives the statement and drop-in text.

% ------------------------------------------------------------------------------------------
\subsection{(a) The adiabatic prescription}\label{pr:a}

\paragraph{a.1 Where on-shell denominators occur.} In light-cone perturbation theory with the valence charges static ($[H_0,\rho]=0$ in the eikonal Hamiltonian, Eq.~(\dref{eq:Hdecomp})), the adiabatically switched evolution $U_\varepsilon(0,\mp\infty)=T\exp\{-i\int_{\mp\infty}^0dx^+e^{-\varepsilon|x^+|}V_I(x^+)\}$ acting on an eigenstate $\ket m$ of $H_0$ gives the Gell-Mann--Low states
\begin{align}
\ket{m,\pm}=\ket m+\sum_n\ket n\frac{V_{nm}}{E_m-E_n\pm i\varepsilon}+\sum_{n,l}\ket n\frac{V_{nl}V_{lm}}{(E_m-E_n\pm2i\varepsilon)(E_m-E_l\pm i\varepsilon)}+\dots
\end{align}
(first order: $-i\int_{-\infty}^0dt\,e^{\varepsilon t}e^{i(E_n-E_m)t}=1/(E_m-E_n+i\varepsilon)$), with the vertex acting later to the left, as in Eq.~(\dref{eq:lcpt}); $+$ is the in-state (prepared at $x^+\to-\infty$), $-$ the out-state. In the spectrum, Eq.~(\dref{eq:observable}), $\Omega$ acts on the soft vacuum: $E_0-E_n<0$ for every soft Fock state, so $\mf N$, $\mf A^{(1)}$, $\mf A^{(3)}_2$, $\mf B_2$ are prescription independent. $\bar\Omega^\dagger$ is the final-state dressing of the observed gluon, and its coefficients follow from $\Omega\ket{g(k)}$: $\mf C_1$ and the $\Theta_<$ branch of $\mf B_3$ have $E_g(k)-E_{gg}\le0$, zero only at the collinear point (the boundary of a positive quadratic form, irrelevant in dimensional regularisation); the only denominator that vanishes \emph{inside} the integration region is $E_g(k)-E_g(p)=\frac{p^+\v k^2-k^+\v p^2}{2k^+p^+}$ of $\mf B_3$ (instantaneous term, the $p^+>k^+$ and $p^+<k^+$ triple-gluon terms, and the commutator part of the two-valence-vertex terms; the anticommutator part has no pole, since $\frac{\rho^b\rho^a}{E_k}-\frac{\rho^a\rho^b}{E_p}$ over $E_k-E_p$ has residue $\propto[\rho^b,\rho^a]$). In the cross section $\mf B_3$ appears only in $\mathcal A_{B_3}$ of Eq.~(\dref{eq:A3amp}) (including its unitarity image $-\bar{\mf B}_3^\dagger\bar{\mf A}$ in $\mf A^{(3)}_1$), i.e.\ only in Row~I.III and its complex conjugate. No other row has an on-shell denominator (Group~II Row~I has coordinate-space denominators vanishing on surfaces, but they are not energy denominators and the sum is continuous, Section~\dref{sec:checks}). Since the observed gluon is an asymptotic final-state gluon, the time-dependent definition uses $\Omega_+$ for the initial state and $\Omega_-$ for the final-state dressing: $\mf B_3$ with $E_k-E_p-i\varepsilon$. Note that the referee's ``in-state'' prescription is relevant only through the final-state dressing; the in-state itself has no vanishing denominators at this order. The PV prescription of the draft is, at this order, $\frac12(\Omega_++\Omega_-)$ (standing waves).

\paragraph{a.2 The on-shell part is anti-Hermitian.} Write the pole part of $\mf B^{ba}_{3ji}(k\to p)$, Eq.~(\dref{eq:B3mom}), as $ig^2f^{abc}\rho^c(\v k-\v p)\,n_{ji}(k,p)/D$, $D=p^+\v k^2-k^+\v p^2$. Then $1/(D\mp i0)=\mathcal P\frac1D\pm i\pi\delta(D)$ adds $\Delta_\mp=\pm i\pi\,ig^2f\rho\,n\,\delta(D)$. $\Delta$ is anti-Hermitian iff $n_{ji}(k,p)=n_{ij}(p,k)$ on the shell. \emph{Check} (a1): over 2000 random on-shell points, both branches ($p^+<k^+$ maps to $p^+>k^+$ under $k\leftrightarrow p$), maximal relative violation $3.8\times10^{-13}$; the double pole $1/(p^+-k^+)^2$ cancels on shell between the instantaneous and the triple-gluon term ($n$ finite at $p^+\to k^+$), which fixes the relative sign that the symmetry test alone does not. Consequences: $\Omega_\pm=\Omega_{\rm PV}+\Delta_\pm$ with $\Delta_\pm^\dagger=-\Delta_\pm$ at this order; the $a^\dagger a$ relation of $\Omega^\dagger\Omega=1$ (Appendix~\dref{app:unitarity}) contains only $\mf B_3+\mf B_3^\dagger$ and is satisfied for all three prescriptions, so unitarity cannot fix the prescription. In the amplitude $\mathcal A_{B_3}=\int\bar{\mf B}_3^\dagger[U\mf A-\bar{\mf A}]$ the bracket is the LO amplitude of the integrated gluon, so
\begin{align}\label{pr:DNgen}
\Delta N(k)=\big\langle\mathcal A_{\rm LO}^\dagger(k)\,\bar\Delta^\dagger\mathcal A_{\rm LO}\big\rangle_k+{\rm h.c.}=\int_p\big[G(k,p)-G(p,k)\big],\qquad G(k,p)\equiv\big\langle\mathcal A_{\rm LO}^\dagger(k)\,\bar\Delta^\dagger_{kp}\,\mathcal A_{\rm LO}(p)\big\rangle\,,
\end{align}
antisymmetric in $k\leftrightarrow p$: the total number of gluons is unchanged, $\int d^3k\,\Delta N=0$, but the spectrum at fixed $k$ is changed by on-shell rescattering of the produced gluon on the (rotated) valence charges after the collision. It is linear in $\eta$: in- and out-dressing give opposite signs.

\paragraph{a.3 From momentum to coordinate space.} $E_k-E_p\pm i\varepsilon$ is $p^+\to p^+\pm i0$ ($\partial_{p^+}D=\v k^2>0$). The coordinate forms are boundary values of functions analytic in $\xi=p^+/k^+$: for ${\rm Im}\,\xi>0$, with $1/X=-i\int_0^\infty d\tau\,e^{i\tau X}$ (convergent, ${\rm Im}X>0$), both Gaussian integrals converge, and
\begin{align}
\int d^2k\,d^2p\,\frac{e^{-i\v k\cdot\v a+i\v p\cdot\v b}}{\xi\v k^2-\v p^2}=-\frac{i\pi^2}{\xi}\int_0^\infty ds\,e^{\frac{is}{4}(\v b^2-\v a^2/\xi)}=\frac{4\pi^2}{\xi\v b^2-\v a^2}\qquad({\rm Im}\,\xi>0)\,,
\end{align}
whose boundary value at $\xi+i0$ is $4\pi^2/(\xi\v b^2-\v a^2+i0)$. Hence $D\pm i0$ in momentum space is $D_x\pm i0$ for the coordinate denominator $D_x=p^+(\v x-\v u)^2-k^+(\v x-\v z)^2$ of Eq.~(\dref{eq:B3}) ($\v u$ conjugate to $\v p$). In Row~I.III, which contains $\bar{\mf B}_3^\dagger$, the out-state dressing gives $1/(D+i0)$ with $D\equiv p^+(\v x-\v z)^2-k^+(\v x-\v w)^2$ ($\v z$: integrated gluon, $\v w$: observed gluon); we write $1/(D+i\eta0)$, $\eta=+1$ (out), $-1$ (in), and PV for $\eta=0$.

\paragraph{a.4 Row I.III: the residues.} Before the $p^+$ integral (Group\_I\_checked, Row~III), in units of $\frac{1}{(2\pi)^3}\frac{ig^4f^{abd}}{8\pi^4}\frac{dp^+}{2\pi}$ times the colour factor $\mathcal X\equiv\mathcal L^a\big[U^{de}(x)U^{bc}(z)-U^{de}(x)U^{bc}(y)\big]\rho^e(x)\rho^c(y)$, with $a=x-w$, $b=x-z$, $c=z-w$, $Y=y-z$, $A=a^2$, $B=b^2$, $P=\Kv(x'-w')$, the five terms are
\begin{align}
&T_1=\frac{P\cdot Y}{Y^2}\frac{p^++k^+}{(k^+-p^+)^2D}\,,\qquad T_3=-\frac{(P\cdot a)(b\cdot Y)}{ABY^2}\frac{k^+A+p^+B}{p^+k^+D}\qquad(p^+\in(\Lambda,k^+-\Lambda)\cup(k^++\Lambda,\vee))\,,\notag\\
&T_4=\frac{2}{k^+}\frac{\mathcal T[V']}{Y^2c^2(k^+-p^+)D}\ \ (p^+>k^+)\,,\qquad T_5=\frac{2p^+\,\mathcal T[V']}{Y^2|V'|^2(k^+-p^+)D}\ \ (p^+<k^+)\,,\qquad V'=k^+a-p^+b\,,\notag\\
&\mathcal T[V']=(P\cdot c)(Y\cdot V')-\frac{k^+}{p^+}(Y\cdot c)(P\cdot V')-\frac{k^+}{k^+-p^+}(P\cdot Y)(c\cdot V')\,.
\end{align}
$1/(D+i\eta0)=\mathcal P\frac1D-i\eta\pi\delta(D)$, $\delta(D)=\delta(p^+-p_0)/B$, $p_0=k^+A/B$. At $p_0$, $|V'|^2=k^{+2}\frac AB c^2$, so $T_4$ and $T_5$ have the same residue and one expression holds on both sides of $A=B$:
\begin{align}\label{pr:R}
R&=\sum_n\frac{[DT_n]_{p_0}}{B}=\frac{1}{k^+}P^i\Bigg\{\frac{Y^i(A+B)}{Y^2(B-A)^2}-\frac{2a^i(b\cdot Y)}{ABY^2}+\frac{2\big[c^i(Y\cdot V)-\frac BA(Y\cdot c)V^i-\frac{B}{B-A}Y^i(c\cdot V)\big]}{Y^2c^2(B-A)}\Bigg\},\quad V=a-\tfrac ABb\,.
\end{align}
\emph{Check} (a2): $R$ against the direct $p^+$ integral of $\sum_nDT_n$ with the Lorentzian $\frac1\pi\frac{\varepsilon}{D^2+\varepsilon^2}$ over the exact ranges ($\Lambda/k^+=10^{-7}$, $\vee/k^+=10^7$), eight random configurations with $A/B$ from $0.14$ to $10.5$, both branches: e.g.\ $-0.6608715$ against $-0.66087283$ ($\varepsilon=10^{-7}$), $-0.4584715$ against $-0.4584714$, $0.00215016$ against $0.00215014$; the approach is $O(\varepsilon)$ in all cases (two cases with a large $O(\varepsilon)$ coefficient converge more slowly). Also at $|\v z-\v w|=10^{-2}$: $2161.698$ against $2161.861$ ($\varepsilon=10^{-7}$), and at $|\v z|=20$: $-7.467\times10^{-5}$ against $-7.480\times10^{-5}$ (a5).

\paragraph{a.5 Exact identity: the $i\pi\theta$ terms become a constant $i\pi$ in the combined kernel.} Writing $R=P^iQ^{ij}\Kv^j(y-z)$, sympy gives exactly (a3)
\begin{align}\label{pr:QB}
Q^{ij}&=\frac{2}{k^+}\,\mathfrak B^{ij}\,,\qquad
\mathfrak B^{ij}\equiv\frac{\delta^{ij}}{2(z-w)^2}-\frac{(x-w)^i(x-z)^j}{(x-w)^2(x-z)^2}\notag\\&\qquad+\frac{(x-w)^i(z-w)^j\frac{(x-z)^2}{(x-w)^2}+(z-w)^i(x-z)^j\frac{(x-w)^2}{(x-z)^2}-(z-w)^i(x-w)^j-(x-z)^i(z-w)^j}{(z-w)^2[(x-w)^2-(x-z)^2]}\notag\\
&=\frac{\cos2\phi\,\delta^{ij}+\sin2\phi\,\epsilon^{ij}}{2(\v z-\v w)^2}\,,\qquad \cos2\phi=\frac{[a\cdot b]^2-[a\times b]^2}{AB}\,,\ \ \sin2\phi=\frac{2(a\cdot b)(a\times b)}{AB}\,,\ \ \epsilon^{12}=+1\,,
\end{align}
where $\mathfrak B^{ij}$ is \emph{literally the brace multiplying} $\ln\frac{(x-z)^2}{(x-w)^2}$ in $\mathcal M^{ij}$, Eq.~(\dref{eq:Mij}), and $\phi$ the angle from $\v x-\v w$ to $\v x-\v z$ (numerically, the compact form agrees with the residue sum~\eqref{pr:R} at $2\times10^4$ random points; it is manifestly regular on the line $(\v x-\v z)^2=(\v x-\v w)^2$, where the separate residues have double poles). With the prefactors, the extra term of Row~I.III is $2i\mathcal C_{\rm f}f^{abd}\int P^i\Kv^j(y-z)(-i\eta\pi)\mathfrak B^{ij}\mathcal X$, i.e.
\begin{align}\label{pr:shift}
\boxed{\ \mathcal M^{ij}\Big|_{i\eta0}=\mathcal M^{ij}\Big|_{\rm PV}-i\eta\pi\,\mathfrak B^{ij}\,:\qquad \ln\frac{(x-z)^2}{(x-w)^2}\ \to\ \ln\frac{(x-z)^2}{(x-w)^2}-i\eta\pi\ \text{ in Eq.~(\dref{eq:Mij})}\ }
\end{align}
The other logarithms of $\mathcal M^{ij}$ and $\mathcal A_0$ are unchanged (no pole), as is the term $(\ell-\mathcal E)$ of Eq.~(\dref{eq:GIIII}). So the $i\pi\theta\big((x-z)^2-(x-w)^2\big)$ and $i\pi\theta\big((x-w)^2-(x-z)^2\big)$ pieces of $N_1,N_2,N_{3a}$ do not cancel in the combined kernel: they add up to $i\pi$ on both sides of the line. The pole on the line cancels for any $\eta$.

\paragraph{a.6 Integrability; the forward Coulomb logarithm and its cancellation.} Limits of $R$ (a5): bounded at $\v z\to\v x$ and $\v w\to\v x$, $\propto|\v z-\v y|^{-1}$ at $\v z\to\v y$, $O(|\v z|^{-3})$ at large $\v z$ (no $1/\v z^2$ tail in any direction), $O(|\v w|^{-2})$ at large $\v w$ (Fourier convergent). The single non-integrable point is $\v z\to\v w$, where $\phi\to0$ and
\begin{align}
R\to\frac{1}{k^+}\frac{P\cdot\Kv(y-w)}{(\v z-\v w)^2}\qquad\text{(angular mean to }10^{-6}\text{ at }|\v z-\v w|=10^{-4}\text{; anisotropy }O(|\v z-\v w|)).
\end{align}
Here $p_0\to k^+$, and the only regulator is the window $|p^+-k^+|>\Lambda$ of $\mf B_3$: $|A-B|>\frac{\Lambda}{k^+}B$, near $\v w$ a band of half-width $\frac{\Lambda}{2k^+}|\v x-\v w|$ around the circle $|\v x-\v z|=|\v x-\v w|$. Then (a4: $2\pi\ln(\rho_0/\lambda)+2\lambda/\rho_0$ numerically for $\lambda=10^{-2},10^{-3},10^{-4}$)
\begin{align}
\int_{|\v z-\v w|<\rho_0}\frac{d^2z\ \theta_{\rm window}}{(\v z-\v w)^2}=2\pi\ln\frac{\rho_0k^+}{\Lambda|\v x-\v w|}+O\Big(\frac{\Lambda}{k^+}\Big),
\end{align}
a logarithm $\ln(k^+/\Lambda)$ multiplying, at $\v z=\v w$, $\mathcal X(w)=f^{abd}\mathcal L^a\,\bar Q^d\,[U(w)-U(y)]^{bc}\rho^c(y)$ with the kernel $P\cdot\Kv(y-w)$ and no kernel in $\v x$: $\bar Q^d=\int_xU^{de}(x)\rho^e(x)$ is the total colour charge of the projectile after the collision, and the structure is $\mathcal A_{\rm LO}^{a\dagger}\,f^{abd}\bar Q^d\,\mathcal A_{\rm LO}^b$: the forward (Coulomb) scattering phase of the produced gluon. Under the relabelling $(\v x',\v w')\leftrightarrow(\v y,\v w)$, $b'\leftrightarrow c$, $a\leftrightarrow b$ (which keeps the order of the charges, $P\cdot\Kv(y-w)$ and $\bar Q$ in the middle) the Hermitian conjugate of this term is exactly minus the term (operator identity). Hence the $\ln(k^+/\Lambda)$ cancels between Row~I.III and its complex conjugate, as a Coulomb phase must, and leaves $\ln\frac{|\v x-\v w'|}{|\v x-\v w|}$. Cut-offs at $p_0=\Lambda$ (large $\v z$) and $p_0=\vee$ ($\v z\to\v x$) cost $O(\sqrt{\Lambda/k^+})$ and $O(k^+/\vee)$.

\paragraph{a.7 Result.} With $\mathcal X(z)$ as above, $\mathcal X(w)$ its value at $\v z=\v w$, and any $\rho_0>0$ (the result does not depend on it),
\begin{align}\label{pr:DN}
\frac{d^3\Delta N^{\rm I.III}}{d^3k}&=2\eta\pi\,\mathcal C_{\rm f}\,f^{abd}\int_{w,w'}e^{-ik\cdot(w-w')}\int_{x,y,x'}\Bigg\{\int_z\Big[P^i\mathfrak B^{ij}(x,z,w)\,\Kv^j(y-z)\,\mathcal X(z)-\theta(\rho_0-|\v z-\v w|)\frac{P\cdot\Kv(y-w)}{2(\v z-\v w)^2}\,\mathcal X(w)\Big]\notag\\
&\hspace{6.2cm}+\pi\ln\frac{\rho_0}{|\v x-\v w|}\,P\cdot\Kv(y-w)\,\mathcal X(w)\Bigg\}+{\rm h.c.}
\end{align}
($\eta=+1$ out-state, the physical final-state dressing). Properties: (i) every term is integrable [$\mathcal X(z)-\mathcal X(w)=O(|\v z-\v w|^\alpha)$ for H\"older-$\alpha$ Wilson lines]; (ii) it does not cancel with the complex conjugate: the PV terms of Row~I.III are $i\times$(real kernel), $\Delta N$ is (real kernel)$+$h.c.; in the language of Section~\dref{sec:fin_conv}, its symmetrised three-$\rho$ part is \emph{even} in $\v k$ and its commutator (two-$\rho$) part \emph{odd}, the opposite of the PV terms; (iii) $\int d^3k\,\Delta N=0$ (Eq.~\eqref{pr:DNgen}); (iv) after an average over an ensemble that is invariant under rotations by $\pi$ the odd part vanishes, so $\Delta N$ contributes to the averaged spectrum only through the symmetrised three-point function $\langle\rho\rho\rho\rangle_{\rm sym}$, which vanishes for a Gaussian projectile; event by event (and in correlations) it contributes through the commutator terms (symmetry argument, not an explicit average); (v) the size of $\Delta N$ for a realistic configuration was \emph{not} evaluated. Since $\mf B_3$ enters no other row and no other coefficient has on-shell denominators, nothing else can cancel $\Delta N$.

\paragraph{a.8 Answer to Major 5, first point.} The $i\pi\theta$ terms do not cancel in the combined kernel; there they become $-i\eta\pi$ times the coefficient of $\ln\frac{(x-z)^2}{(x-w)^2}$, Eq.~\eqref{pr:shift}. With the complex conjugate they leave the finite term~\eqref{pr:DN}. The pole cancellation and all integrability statements of Section~\dref{sec:finite} hold for any $\eta$. \textbf{Correction needed:} the parenthetical sentence at the end of Section~\dref{sec:fin_pv} (``would add terms $\mp i\pi\theta$ \dots\ to the separate terms'') understates the effect, and condition (C6) must say that the result is the spectrum of standing-wave (PV) dressed gluons; for out-state dressed gluons $\Delta N$ has to be added. Physically, the time-dependent definition of an observed final-state gluon is the out-state, so I recommend adding $\Delta N$ with $\eta=+1$ to the finite terms (or, at least, giving it as the prescription dependence). Caveat: with static charges the final-state interaction lasts for ever; whether the full on-shell part is physical beyond the eikonal approximation, where the valence partons and the gluon separate, is a question this calculation cannot answer.

% ------------------------------------------------------------------------------------------
\subsection{(b) Instantaneous Group III terms}\label{pr:b}

\paragraph{b.1 Where in $p^+$ the finite part comes from.} After the $\v w'$ integral, the terms $-\frac{1}{k^+}\delta\delta$ of Rows~III.III, III.V, III.VI are, with $t=p^+/k^+$, $\v r=\v y'-\v z$, $\v q=(1+t)\v k$ and $b(\v r)=\Kv^j(x'-y')\mathcal L(y')|_{y'=z+r}$,
\begin{align}
I(T)=\int_0^Tdt\,h(t)\,,\qquad h(t)=\int d^2r\,e^{i\v q\cdot\v r}\,\Kv^m(\v r)\big[b(\v r)-b(0)\big]\,,\qquad T=\frac{\vee}{k^+}\,.
\end{align}
Pointwise in $\v r$ the integrand does not decay in $t$ (the referee's point); after the $\v r$ integral it does. $\Kv^m(\v r)[b(\v r)-b(0)]\to r^mr^n\partial_nb(0)/\v r^2$ is homogeneous of degree zero, so for a $b$ that is differentiable at $\v r=0$, $h(t)\to-2\pi(2\hat q^m\hat q^n-\delta^{mn})\partial_nb(0)/\v q^2\propto t^{-2}$; for a H\"older-$\alpha$ bracket (Wilson lines of an MV-type target), $h\propto t^{-1-\alpha}$. Large $p^+$ corresponds to $|\v r|\sim1/(t|\v k|)$: the hard daughter close to its parent, the observed gluon at $\sim1/|\v k|$. Model (b1): $b(\v r)-b(0)=(\v u\cdot\v r)|\v r|^{\alpha-1}e^{-\v r^2/2s^2}$, closed-form $h$:
\begin{center}\footnotesize
\begin{tabular}{lcccccc}
fraction of $\int_0^\infty h$ from $p^+<t_0k^+$, $t_0=$ & 1 & 3 & 10 & 30 & 100 & 1000\\\hline
$\alpha=1$, $s=1/|\v k|$, component $\parallel\v k$ & $-0.10$ & $0.36$ & $0.77$ & $0.92$ & $0.975$ & $0.9975$\\
$\alpha=1$, $s=1/|\v k|$, component $\perp\v k$ & $0.38$ & $0.68$ & $0.89$ & $0.96$ & $0.988$ & $0.9987$\\
$\alpha=\frac12$ (rough), $s=1/|\v k|$, $\parallel$ & $-0.30$ & $-0.08$ & $0.33$ & $0.61$ & $0.78$ & $0.93$\\
$\alpha=1$, $s=0.1/|\v k|$, $\perp$ & $0.04$ & $0.12$ & $0.37$ & $0.74$ & $0.92$ & $0.992$
\end{tabular}\end{center}
$t^2h(t)\to-6.275,\ 3.138$ at $t=1600$ against $-2\pi u_\parallel$, $+2\pi u_\perp$; $t^{3/2}h\to\mp3.000$ for $\alpha=\frac12$. So the finite part comes from $p^+\lesssim k^+\max\{1,1/(|\v k|\ell)\}$, with $\ell$ the shortest scale on which the bracket varies ($1/Q_s$, the distance $|\v x'-\v z|$); $p^+>P$ contributes $\simeq2\pi|\partial b(0)|\,k^+/(\v k^2P)$ in units of $k^+$, i.e.\ $O(k^+/P)$, or $O((k^+/P)^\alpha)$. The eikonal requirement is $\vee\gg k^+\max\{1,1/(|\v k|\ell)\}$, i.e.\ for $|\v k|\ll Q_s$: $\ln\frac{\vee}{k^+}\gg\ln\frac{Q_s}{|\v k|}$; nothing comes from $p^+\sim\vee$ beyond a power correction.

\paragraph{b.2 Smooth separation of soft and valence modes.} Replace $\theta(\vee-p^+)$ by $f(p^+/\vee)$, $f(0)=1$, $f\to0$ fast enough. Since $h\in L^1$, $\int_0^\infty dt\,f(t/T)h(t)\to\int_0^\infty h$ by dominated convergence, for every $f$; the deviation is $O(1/T)$ for $f'(0)=0$ and $O(\ln T/T)$ otherwise. Numerically (b1, $\parallel$ component, deviation from the limit at $T=10,100,1000$): sharp $0.57,\,0.062,\,0.0059$; $e^{-u}$: $1.10,\,0.24,\,0.038$; $e^{-u^2}$: $0.83,\,0.105,\,0.0106$; $1/(1+u^2)$: $0.72,\,0.093,\,0.0094$. The subtraction term of the bracket (the BFKL piece, $\propto\int dt\,f(t/T)/(1+t)$) gives $\ln T+c_f+O(1/T)$ with $c_f=\int_0^\infty\frac{du}{u}[f(u)-\theta(1-u)]$ (numerically $0,-\gamma_E,-\gamma_E/2,0$ for the four $f$, to $10^{-4}$ at $T=10^4$). In every row the $\ln\vee$ arises from $\int^\vee dp^+/p^+$ with integrands $\propto1/p^+$ up to $O(k^+/p^+)$, so the same $c_f$ appears in all rows: the whole $O(\alpha_s)$ result depends on $f$ only through $\ln\frac{\vee}{k^+}\to\ln\frac{\vee}{k^+}+c_f$ in the BFKL term, which is the change of the rapidity-factorisation scale $\vee\to e^{c_f}\vee$ compensated by the evolution of the valence distribution; the finite parts are $f$-independent up to $O(k^+/\vee)$. At finite $\vee$ every integral is absolutely convergent, so condition (C7) on the order of integration is not an extra assumption: the distributional form~(\dref{eq:instdistr}) is the $\vee\to\infty$ limit of the cut-off integral, for any $f$.

\paragraph{b.3 Answer to Major 5, second point.} The finite parts of the instantaneous terms are under control: they do not come from $p^+\sim\vee$, they are independent of the shape of the soft/valence separation, and the approach to the limit is a power law. Caveat: the transverse test functions are models with the correct short-distance behaviour; the actual bracket with Wilson lines was not integrated.

% ------------------------------------------------------------------------------------------
\subsection{(c) The leftover pole $\propto\ell_\vee$ and the order of limits}\label{pr:c}

\paragraph{c.1 Origin.} In the one-$\rho$ part of $\mf A^{(3)}$ (Appendix~\dref{app:wavefunctions}, ``Sum of the one-$\rho$ terms'') the parent of the observed gluon carries $k^+/\xi$ with $\frac{k^+}{\vee}<\xi<\frac{k^+}{k^++\Lambda}$, and $\int d\xi\,\frac{2-\xi}{\xi(1-\xi)}=\int d\xi\big[\frac2\xi+\frac1{1-\xi}\big]$. The $\frac2\xi$ part (enhanced at small $\xi$, the BFKL region) gives $2\ln\frac\vee{k^+}$, which cancels in the pole coefficient against the $\ln\vee$ of the other diagrams; the $\frac{1}{1-\xi}$ part gives $L+\ell_\vee$ with
\begin{align}
\ell_\vee=\ln\Big(1-\frac{k^+}{\vee}\Big)=-\int_0^{k^+/\vee}\frac{d\xi}{1-\xi}\,,
\end{align}
the part of the \emph{non-enhanced} integrand that is missing because parents faster than $\vee$ (valence modes) are excluded. The infrared $-\ell_\vee$ of Row~II.IV has the same origin (range $p^++k^+<\vee$). With a smooth separation $f$ (c1): $\ell_\vee\to\ell_f=-\int_0^1d\xi\frac{1-f(k^+/\xi\vee)}{1-\xi}=-\kappa_f\frac{k^+}{\vee}+o\big(\frac{k^+}{\vee}\big)$, $\kappa_f=\int_0^\infty dv\frac{1-f(v)}{v^2}$; numerically $\ell_f/(k^+/\vee)=-1.00005,\,-1.777,\,-1.575$ at $k^+/\vee=10^{-4}$ against $\kappa_f=1,\sqrt\pi,\pi/2$ (sharp, $e^{-u^2}$, $1/(1+u^2)$); for $f'(0)\neq0$ the observed gluon itself gets the weight $f(k^+/\vee)$ and $\ell_f=O(\frac{k^+}{\vee}L)$. The BFKL logarithm shifts only by the universal $c_f$.

\paragraph{c.2 Statement.} (1) The remainder is $O(k^+/\vee)\times\frac1\epsilon$: $-\frac{\ell_\vee}{2}\frac2\epsilon=\frac{k^+}{2\vee}\frac2\epsilon+\dots$; it vanishes for $\vee\to\infty$ at fixed $\epsilon$ and diverges for $\epsilon\to0$ at fixed $\vee$. (2) Its coefficient depends on the shape of the soft/valence separation, so it is not a quantity of the effective theory. (3) It comes from parent momenta above $\vee$ in a part of the integrand that the static source does not describe (the source $\rho$ represents the valence modes correctly only for the emission of gluons much softer than themselves, the $1/\xi$-enhanced part; the non-enhanced part is of the order of the neglected recoil terms of the valence partons, $k^+/P^+_{\rm valence}$). (4) The pole factor is a product of a longitudinal and a transverse factor, $\big(\frac{11}{3}-2L+\ell_\vee\big)\big[-\frac2\epsilon+\ln\frac{\v k^2}{\mu^2}\big]$, so the expansion in $k^+/\vee$ is uniform in the transverse loop momentum and commutes with the transverse integral at any fixed $\epsilon$. (5) A UV pole of a power-suppressed term would have to be renormalised, at that power, by non-eikonal couplings of the valence sector, which are not in the Hamiltonian~(\dref{eq:Hdecomp}); at leading power it is absent. The same holds for the $O(\Lambda/k^+)$ terms on the target side. Hence the rule is: \emph{the eikonal approximation is the leading term of the expansion in $k^+/\vee$ and $\Lambda/k^+$ (with logarithms kept), taken at the level of the longitudinal integrands, at fixed transverse momenta and fixed $\epsilon$, before the transverse integrals are done and before $\epsilon\to0$.} It is not a cancellation.

% ------------------------------------------------------------------------------------------
\subsection{Drop-in text for the paper}\label{pr:dropin}
% In the boxes the labels are those of the draft (\eqref, \ref); in this note they are printed as their names.
\begingroup\renewcommand{\eqref}[1]{(\dref{#1})}\renewcommand{\ref}[1]{\dref{#1}}

\paragraph{Section 6.4 (\dref{sec:fin_pv}), replace the paragraph from ``In $\mf B_3$ the denominators'' to the end by:}\mbox{}\\[2pt]
\noindent\fbox{\parbox{0.97\linewidth}{\small
In $\mf B_3$ the denominators $p^+\v k^2-k^+\v p^2$ (momentum space) and $p^+(\v x-\v u)^2-k^+(\v x-\v z)^2$ (coordinate space) are, up to positive factors, the differences $E_g(k)-E_g(p)$ of the light-cone energies of two one-gluon states connected by the static charge, which transfers transverse and longitudinal momentum but no light-cone energy; they vanish on the energy shell. They are the only vanishing energy denominators of the calculation: the vacuum dressing has $E_0-E_n<0$, and the splitting denominators $E_g-E_{gg}$ vanish only at the collinear point. They enter the cross section only through the final-state dressing $\bar\Omega^\dagger$, i.e.\ only in Row~I.III and its complex conjugate. We take them as principal values, which corresponds to standing-wave dressed gluons. In the adiabatic (time-dependent) construction the observed gluon is an out-state, $E_g(k)-E_g(p)-i\varepsilon$ in $\mf B_3$, i.e.\ $1/(D+i\eta0)$ with $\eta=1$ for $D=p^+(\v x-\v z)^2-k^+(\v x-\v w)^2$ in Row~I.III ($\eta=-1$ for an in-state dressing of the final state). With $A=(\v x-\v w)^2$, $B=(\v x-\v z)^2$, $\int_{p_1}^{p_2}dp^+\frac{1}{p^+B-k^+A+i\eta0}=\frac1B\big[\ln\big|\frac{p_2B-k^+A}{p_1B-k^+A}\big|-i\eta\pi\,\theta(p_1<k^+A/B<p_2)\big]$. In the separate terms $N_1$, $N_2$, $N_{3a}$ every logarithm whose argument changes sign is a logarithm of the absolute value,
$\ln\big|\frac{(x-z)^2}{(x-z)^2-(x-w)^2}\big|$ ($N_1,N_2$) and $\ln\big|\frac{(x-w)^2-(x-z)^2}{(x-w)^2}\big|$ ($N_1,N_{3a}$), accompanied by $-i\eta\pi\,\theta$ terms; the other logarithms have positive arguments, and $\mathcal A_0$ is real. In the combined kernel~\eqref{eq:Mij} these logarithms combine into $\ln\frac{(x-z)^2}{(x-w)^2}$, the pole on the line $(\v x-\v z)^2=(\v x-\v w)^2$ cancels for every $\eta$, and the $\theta$ terms combine into a constant: exactly,
\begin{align*}
\mathcal M^{ij}\big|_{i\eta0}=\mathcal M^{ij}\big|_{\rm PV}-i\eta\pi\,\mathfrak B^{ij},\qquad \mathfrak B^{ij}=\frac{\cos2\phi\,\delta^{ij}+\sin2\phi\,\epsilon^{ij}}{2(\v z-\v w)^2},
\end{align*}
where $\mathfrak B^{ij}$ is the brace that multiplies $\ln\frac{(x-z)^2}{(x-w)^2}$ in Eq.~\eqref{eq:Mij} and $\phi$ the angle from $\v x-\v w$ to $\v x-\v z$. The on-shell part of $\mf B_3$ is anti-Hermitian (a unitary rotation within the energy shell, compatible with $\Omega^\dagger\Omega=1$), so the difference conserves the total number of gluons, but it changes the spectrum: since the principal-value terms of Row~I.III are $i$ times a real kernel and the new term is real, the complex conjugate does not cancel it. The difference between out-state and standing-wave dressing is
\begin{align*}
\frac{d^3\Delta N}{d^3k}&=2\eta\pi\,\mathcal C_{\rm f}\,f^{abd}\!\int_{w,w'}\!e^{-ik\cdot(w-w')}\!\int_{x,y,x'}\!\Big\{\!\int_z\!\Big[P^i\mathfrak B^{ij}\Kv^j(y-z)\,\mathcal X(z)-\theta(\rho_0-|\v z-\v w|)\frac{P\cdot\Kv(y-w)}{2(\v z-\v w)^2}\mathcal X(w)\Big]\\&\hspace{6cm}+\pi\ln\frac{\rho_0}{|\v x-\v w|}\,P\cdot\Kv(y-w)\,\mathcal X(w)\Big\}+{\rm h.c.},
\end{align*}
with $\mathcal X(z)=\mathcal L^a[U^{de}(x)U^{bc}(z)-U^{de}(x)U^{bc}(y)]\rho^e(x)\rho^c(y)$ and any $\rho_0$. The non-integrable $1/(\v z-\v w)^2$ of the first term is cut off by the window $|p^+-k^+|>\Lambda$ into $2\pi\ln\frac{\rho_0k^+}{\Lambda|\v x-\v w|}$; the $\ln(k^+/\Lambda)$ multiplies $\mathcal A^\dagger_{\rm LO}f\cdot\bar Q\,\mathcal A_{\rm LO}$, $\bar Q=\int_xU(x)\rho(x)$, a forward Coulomb phase, and cancels exactly against the complex conjugate. $\Delta N$ is finite and integrable; its symmetrised three-$\rho$ part is even and its commutator part odd in $\v k$. The finite terms given in this paper are those of the principal-value prescription; for out-state dressed gluons $\Delta N$ (with $\eta=1$) is to be added.
}}

\paragraph{Condition (C6), replace by:} \emph{(C6) The vanishing energy denominators of $\mf B_3$ are principal values (standing-wave dressing of the final state); the adiabatic out-state dressing adds $\Delta N$ of Section~\ref{sec:fin_pv}, which is finite only together with its complex conjugate.}

\paragraph{Condition (C7), replace by:} \emph{(C7) In the instantaneous terms of Group~III the $p^+$ integral is cut at $\vee$ (or by any smooth $f(p^+/\vee)$); at finite $\vee$ all integrals converge absolutely, and the distributional form~\eqref{eq:instdistr} is the limit $\vee\to\infty$, reached like $k^+/\vee$ (Section~\ref{sec:fin_inst}).}

\paragraph{Section 6.5 (\dref{sec:fin_inst}), append:}\mbox{}\\[2pt]
\noindent\fbox{\parbox{0.97\linewidth}{\small
Although the integrand does not decay in $p^+$ at fixed transverse positions, the terms do not come from $p^+\sim\vee$. After the $\v y'$ integral the $p^+$ distribution is $h(p^+/k^+)$, $h(t)=\int d^2r\,e^{i(1+t)\v k\cdot\v r}\Kv^m(\v r)[b(\v r)-b(0)]$, with $b$ the first term of the bracket at $\v y'=\v z+\v r$; since $\Kv^m(\v r)[b(\v r)-b(0)]$ is of degree zero at $\v r\to0$, $h(t)\simeq-2\pi(2\hat k^m\hat k^n-\delta^{mn})\partial_nb(0)/[(1+t)^2\v k^2]$, so that $p^+>P$ contributes $O(k^+/P)$ [$O((k^+/P)^\alpha)$ if the Wilson lines are only H\"older continuous with exponent $\alpha$]. Large $p^+$ is the region in which the hard daughter is at a distance $k^+/(p^+|\v k|)$ from its parent. The bulk comes from $p^+\lesssim k^+\max\{1,1/(|\v k|\ell)\}$, $\ell$ the shortest transverse scale of the bracket; the eikonal approximation requires $\vee$ to be large compared to this. With a smooth separation $f(p^+/\vee)$ of soft and valence modes the limit is the same for every $f$; the subtraction term, which is part of the BFKL term, becomes $\ln\frac{\vee}{k^+}+c_f$, $c_f=\int_0^\infty\frac{du}{u}[f(u)-\theta(1-u)]$, the same constant as in all other rows, i.e.\ a change of the rapidity-factorisation scale $\vee\to e^{c_f}\vee$.
}}

\paragraph{Section 5.2 (\dref{sec:dg_poles}), ``The balance'', replace from ``The remainder'' to the end of the paragraph by:}\mbox{}\\[2pt]
\noindent\fbox{\parbox{0.97\linewidth}{\small
The remainder $-\frac{\ell_\vee}{2}\frac2\epsilon$ is a pole multiplied by $\ell_\vee=\ln(1-k^+/\vee)=-k^+/\vee+O(k^{+2}/\vee^2)$, a power correction in the ratio of the soft and the valence longitudinal momenta. It is not cancelled, and we do not claim that it is: it is outside the eikonal approximation. $\ell_\vee=-\int_0^{k^+/\vee}d\xi/(1-\xi)$ is the part of the vertex correction (and, with the opposite sign of the pole, of the real splitting) in which the parent of the observed gluon is faster than $\vee$, i.e.\ a valence mode, and which is not enhanced at small momentum fraction; for this part the description of the valence modes by a static source is not valid, and its coefficient is not defined by the effective theory: for a smooth separation $f$ of soft and valence modes $\ell_\vee$ becomes $-\kappa_f k^+/\vee$ with $\kappa_f=\int_0^\infty dv\,[1-f(v)]/v^2$, while $\ln(\vee/k^+)$ only shifts by a universal constant. We therefore define the eikonal approximation as the leading term of the expansion in $k^+/\vee$ and $\Lambda/k^+$, with the logarithms kept, taken at the level of the longitudinal integrands at fixed transverse momenta and at fixed $\epsilon$, i.e.\ before the transverse integrals are done and the regulator is removed. Since the pole factor of Row~I.II is a product of a longitudinal and a transverse factor, this expansion is uniform in the transverse loop momentum; at fixed $\epsilon$ it is equivalent to $\vee\to\infty$ before $\epsilon\to0$. A pole multiplying a power correction would have to be renormalised, at that power, by non-eikonal interactions of the valence sector that are not part of the Hamiltonian~\eqref{eq:Hdecomp}; at leading power it is absent.
}}\par\medskip
\noindent Condition (C1) of Section~\dref{sec:fin_integrable} then reads: \emph{(C1) $\Lambda\to0$ and $\vee\to\infty$ are taken after the subtractions~\eqref{eq:softsplit}, at fixed $\epsilon$ (leading power of the eikonal approximation, Section~\ref{sec:dg_poles}); terms $O(\Lambda/k^+)$ and $O(k^+/\vee)$ are dropped \dots} In Appendix~\dref{app:distributions}, ``Principal values'', add: ``\emph{This is the standing-wave dressing of the final state; the out-state dressing differs by the finite term $\Delta N$ of Section~\ref{sec:fin_pv}.}'' In Section~\dref{sec:lcpt}, after ``understood as principal values'', add the same reference.

\paragraph{Reply to the referee (short).} \emph{Prescription of $\mf B_3$.} We derived the $i\varepsilon$ form from the adiabatic construction. The only vanishing denominators are those of $\mf B_3$, which enter only Row~I.III. The $i\pi\theta$ terms do not cancel in the combined kernel: they combine exactly into $\ln\frac{(x-z)^2}{(x-w)^2}\to\ln\frac{(x-z)^2}{(x-w)^2}\mp i\pi$, and with the complex conjugate they leave a finite term $\Delta N$ (a forward Coulomb logarithm $\ln(k^+/\Lambda)$ that appears at $\v z\to\v w$ cancels against the conjugate). We now state that our result is the spectrum of standing-wave (principal-value) dressed gluons and give $\Delta N$, the difference to out-state dressed gluons, explicitly. \emph{Instantaneous terms.} After the transverse integral their $p^+$ distribution falls like $(k^+/p^+)^2$; the finite part comes from $p^+\lesssim k^+\max(1,1/|\v k|\ell)$, not from $p^+\sim\vee$, and it is independent of the shape of the soft/valence separation, up to the universal shift of $\ln\vee$ in the BFKL term. \emph{Order of limits.} Rephrased as the definition of the eikonal approximation at leading power at fixed $\epsilon$, with the origin and the cut-off dependence of the coefficient $\ell_\vee$.

\endgroup
% ------------------------------------------------------------------------------------------
\subsection{Not done / open}
\noindent (1) $\Delta N$ was not evaluated numerically for a projectile and target configuration (only its kernel, its limits and the cancellation of the Coulomb logarithm, analytically and pointwise); its size relative to the PV finite terms is unknown. (2) Property (iv) of a.7 (vanishing of the averaged $\Delta N$ for a Gaussian projectile in a rotation-invariant ensemble) is a symmetry argument, not an explicit average. (3) Whether the out-state (physical) or the standing-wave prescription is the appropriate definition within the eikonal approximation, where the static charges interact with the produced gluon for an infinite light-cone time, is a physics choice for the authors; the note gives both. (4) The coordinate-space prescription rests on the analyticity argument of a.3 and the residues were checked against the coordinate expressions of Group\_I\_checked; an independent momentum-space evaluation of $\Delta N$ was not done (only the on-shell kernel $n$ in momentum space). (5) (b) uses model transverse functions with the right short-distance behaviour, not the actual brackets with Wilson lines. (6) The claim that every $\ln\vee$ of every row has the same $c_f$ is argued from the $1/p^+$ form of the integrands, not checked row by row.
